\section{Japan's postwar Hiropon crisis and the fate of its survivors}\label{12_philopon-japans-postwar-hiropon-crisis-and-the-fate-of-its-survivors}

\textbf{Most chronic methamphetamine users from Japan's 1945--1956 epidemic simply vanished from the historical record --- not because they recovered, but because nobody systematically tracked them.} The best available evidence suggests that the majority of an estimated two million users stopped when supply collapsed under aggressive law enforcement, but a devastating minority --- perhaps 55,000 to 150,000 individuals --- suffered lasting psychiatric damage ranging from permanent psychosis vulnerability to chronic institutionalization spanning decades. Tatetsu Seijun's landmark 30-year follow-up of nine patients, published in 1997, found persistent ``schizophrenic states'' and ``remarkable disturbance of spontaneity'' three decades after last use. \href{https://webview.isho.jp/journal/detail/abs/10.11477/mf.1405904411}{Isho} This is the longest follow-up study from any methamphetamine epidemic in history, and its findings remain sobering. The absence of systematic cohort tracking constitutes the single largest gap in the evidence base --- what we know comes from a patchwork of hospital case series, police statistics, and the pioneering clinical observations of Japanese psychiatrists who were, quite literally, describing a disease entity the world had never seen before.

\subsection{Two million users in a shattered nation}\label{12_philopon-two-million-users-in-a-shattered-nation}

The epidemic's origins lie in wartime industrial chemistry. Nagai Nagayoshi first synthesized methamphetamine from ephedrine \href{https://bsd.neuroinf.jp/wiki/\%E8\%A6\%9A\%E9\%86\%92\%E5\%89\%A4}{Neuroinf} in 1893, \href{https://www.globalsecurity.org/military/world/japan/history-20c-drugs.htm}{Global Security} and Ogata Akira achieved crystallization in 1919, \href{https://www.sciencedirect.com/science/article/abs/pii/S095539591200196X}{ScienceDirect +2} but it was Dainippon Pharmaceutical Company (大日本製薬, now Sumitomo Pharma) that brought the drug to mass production around 1941 under the brand name \textbf{Philopon} (フィロポン) \href{https://en.wikipedia.org/wiki/Dainippon_Pharmaceuticals}{Wikipedia} --- from the Greek \emph{philo} (love) + \emph{ponos} (labor). \href{https://www.sciencedirect.com/science/article/abs/pii/S095539591200196X}{ScienceDirect} \href{https://flemingsbond.wpcomstaging.com/philopon-a-japanese-murder-drug/}{Fleming's Bond} The Japanese military distributed methamphetamine tablets stamped with the Imperial crest, called \emph{Totsugeki-jō} (突撃錠, ``storming tablets''), to fighter pilots, signal corps, kamikaze crews, and munitions factory workers \href{https://www.vice.com/en/article/how-stigma-created-japans-hidden-drug-problem/}{VICE} as \emph{senryoku zōkyō zai} (戦力増強剤, ``combat-power-enhancing drug''). \href{https://www.zocalopublicsquare.org/the-world-war-ii-wonder-drug-that-never-left-japan/}{Zócalo Public Square} Upon Japan's surrender in August 1945, massive stockpiles in military warehouses, hospitals, and supply caves flooded onto the black market \href{https://jp.drugfreeworld.org/drugfacts/crystalmeth/history-of-methamphetamine.html}{Drugfreeworld} \href{https://www.zocalopublicsquare.org/the-world-war-ii-wonder-drug-that-never-left-japan/}{Zócalo Public Square} and were sold over the counter at neighborhood pharmacies with no legal restriction. \href{https://www.sciencedirect.com/science/article/abs/pii/S095539591200196X}{ScienceDirect} \href{https://www.zocalopublicsquare.org/the-world-war-ii-wonder-drug-that-never-left-japan/}{Zócalo Public Square}

The scale of what followed was unprecedented. By 1954, the National Police Agency estimated \textbf{approximately 550,000 chronic users} (常用者) and \textbf{two million persons with use experience} (元使用者) \href{https://pubmed.ncbi.nlm.nih.gov/2765722/}{PubMed +2} from a total population of roughly 89 million --- meaning one in every 45 Japanese had used methamphetamine. An anonymous Ministry of Welfare survey in May 1954 found \textbf{7.5\% of respondents} acknowledged having tried Hiropon. \href{https://www.zocalopublicsquare.org/the-world-war-ii-wonder-drug-that-never-left-japan/}{Zócalo Public Square} The drug's appeal cut across demographics: veterans habituated during wartime, students cramming for exams, night-shift factory workers, entertainers, and --- increasingly --- juvenile delinquents in the poverty and anomie of occupied Japan. \href{https://www.thefreelibrary.com/Methamphetamine+use+in+Japan+after+the+Second+World+War:...-a0210368292}{The Free Library} \href{https://www.unodc.org/unodc/en/data-and-analysis/bulletin/bulletin_1989-01-01_1_page007.html}{UNODC} A critical turning point came in August 1949, when the Ministry of Health banned tablet and powder production but \textbf{not injectable formulations}, inadvertently shifting the user population toward intravenous administration with far more severe clinical consequences. \href{https://www.unodc.org/unodc/en/data-and-analysis/bulletin/bulletin_1989-01-01_1_page007.html}{unodc}

The legal response escalated through several phases. The \textbf{Stimulant Control Act} (覚醒剤取締法, Act No.~252) was enacted on June 30, 1951, \href{https://ja.wikipedia.org/wiki/\%E3\%83\%A1\%E3\%82\%BF\%E3\%83\%B3\%E3\%83\%95\%E3\%82\%A7\%E3\%82\%BF\%E3\%83\%9F\%E3\%83\%B3}{Wikipedia} criminalizing production, possession, and use with penalties of up to three years' imprisonment. \href{https://www.unodc.org/unodc/en/data-and-analysis/bulletin/bulletin_1989-01-01_1_page007.html}{UNODC +2} Amendments in 1954 raised penalties to five years for first offenses and seven for recidivists, \href{https://www.unodc.org/unodc/en/data-and-analysis/bulletin/bulletin_1989-01-01_1_page007.html}{UNODC} while the Mental Hygiene Law was simultaneously amended to permit compulsory hospitalization (措置入院) of addicts displaying psychotic symptoms. \href{https://www.unodc.org/unodc/en/data-and-analysis/bulletin/bulletin_1989-01-01_1_page007.html}{UNODC} \href{https://www.unodc.org/unodc/en/data-and-analysis/bulletin/bulletin_1989-01-01_1_page007.html}{unodc} A second major amendment in August 1955 banned importation of precursor chemicals (ephedrine, phenylacetone). \href{https://www.npa.go.jp/hakusyo/s55/s550200.html}{National Police Agency +2} Alongside legislation, police centralization in 1953 dramatically improved enforcement capacity, \href{https://www.unodc.org/unodc/en/data-and-analysis/bulletin/bulletin_1989-01-01_1_page007.html}{UNODC +2} and the government established the Stimulant Countermeasures Promotion Headquarters (覚醒剤対策推進本部) in January 1955, \href{https://www.unodc.org/unodc/en/data-and-analysis/bulletin/bulletin_1957-01-01_3_page003.html}{UNODC} launching education campaigns involving \textbf{385,000 posters, 225,000 pamphlets, and hundreds of films and lectures} across 36 local headquarters in 46 prefectures. \href{https://www.unodc.org/unodc/en/data-and-analysis/bulletin/bulletin_1957-01-01_3_page003.html}{UNODC} Arrests peaked at \textbf{55,664 in 1954}, \href{https://www.npa.go.jp/hakusyo/s55/s550200.html}{National Police Agency} then collapsed: \href{https://www.unodc.org/unodc/en/data-and-analysis/bulletin/bulletin_1989-01-01_1_page007.html}{UNODC} \href{https://yakubutsu-bengo.com/statistics/}{Yakubutsu-bengo} 32,143 (1955), approximately 5,000 (1956), and \textbf{just 271 by 1958} \href{https://www.unodc.org/unodc/en/data-and-analysis/bulletin/bulletin_1989-01-01_1_page007.html}{UNODC} \href{https://www.unodc.org/unodc/en/data-and-analysis/bulletin/bulletin_1989-01-01_1_page007.html}{unodc} --- a 99.5\% decline in four years. \href{https://www.unodc.org/unodc/en/data-and-analysis/bulletin/bulletin_1989-01-01_1_page007.html}{UNODC} As historian Miriam Kingsberg argued, Hiropon was ``not embedded in the postwar political economy or culture, making possible its swift suppression.'' \href{https://www.cambridge.org/core/journals/journal-of-asian-studies/article/abs/methamphetamine-solution-drugs-and-the-reconstruction-of-nation-in-postwar-japan/0A0FCE7E74AE7D1ECE9A134114773CE6}{Cambridge Core} Its eradication became a vehicle for national reconstruction, \href{https://www.researchgate.net/publication/282513182_The_forgotten_success_story_Japan_and_the_methamphetamine_problem}{ResearchGate} symbolically restoring Japanese sovereignty and self-discipline. \href{https://www.researchgate.net/publication/259431765_Methamphetamine_Solution_Drugs_and_the_Reconstruction_of_Nation_in_Postwar_Japan}{ResearchGate +2}

\subsection{The clinical landscape Japanese psychiatrists discovered}\label{12_philopon-the-clinical-landscape-japanese-psychiatrists-discovered}

Japanese clinicians at Tokyo Metropolitan Matsuzawa Hospital (松沢病院) --- the epicenter of treatment --- were the first in the world to systematically describe \textbf{methamphetamine psychosis} (覚醒剤精神病) as a distinct clinical entity. \href{https://www.unodc.org/unodc/en/data-and-analysis/bulletin/bulletin_1989-01-01_1_page007.html}{unodc} The condition presented as a paranoid-hallucinatory state virtually indistinguishable from paranoid schizophrenia: \href{https://www.jmedj.co.jp/journal/paper/detail.php?id=4742}{Jmedj +2} persecutory delusions, auditory hallucinations, thought broadcasting, ideas of reference, and pervasive suspiciousness. \href{https://www.jmedj.co.jp/journal/paper/detail.php?id=4742}{Jmedj +2} But beneath the florid psychotic symptoms, clinicians also documented a constellation of ``negative'' features that proved more resistant to treatment: \textbf{anhedonia} (無快感), \textbf{amotivation and loss of spontaneity} (意欲障害/自発性障害), emotional blunting, \href{https://pubmed.ncbi.nlm.nih.gov/11085303/}{PubMed} anxiety, insomnia, and personality deterioration.

Sato Mitsumoto (佐藤光源) of Tohoku University, building on decades of Japanese clinical observation, identified \textbf{three distinct clinical courses} after methamphetamine cessation \href{https://nyaspubs.onlinelibrary.wiley.com/doi/abs/10.1196/annals.1316.035}{Annals of the New York Academy of Sciences} that became the standard classification framework. \href{https://pubmed.ncbi.nlm.nih.gov/15542728/}{PubMed} The \textbf{transient type} saw psychosis resolve quickly, though vulnerability to relapse persisted indefinitely. The \textbf{prolonged type} featured psychosis continuing weeks to months after cessation without re-use. The \textbf{persistent type} constituted a chronic, schizophrenia-like state that became self-sustaining \href{https://pubmed.ncbi.nlm.nih.gov/1553491/}{nih} \href{https://pubmed.ncbi.nlm.nih.gov/1553491/}{PubMed} --- essentially a permanent acquired psychotic illness (Sato M, Numachi Y, Hamamura T. ``Relapse of paranoid psychotic state in methamphetamine model of schizophrenia.'' \emph{Schizophrenia Bulletin} 1992; 18(1): 115--122).

The distinction between acute psychosis and chronic sequelae was documented from the earliest years. Hayashi (1955), studying \textbf{74 patients} in the pre-neuroleptic era, found that approximately two-thirds remitted within 20 days of cessation, but roughly \textbf{10\% displayed psychosis persisting beyond six months} and in some cases years \href{https://psychiatryonline.org/doi/10.1176/appi.ajp.2009.08111695}{Psychiatry Online} (Hayashi S. ``Wake-amine addiction.'' \emph{Sōgō Igaku} 1955; 12: 656--661). Later data from Yui and colleagues refined these figures: \textbf{64\% recovery within 10 days} and \textbf{82\% within 30 days}, \href{https://ja.wikipedia.org/wiki/\%E7\%B2\%BE\%E7\%A5\%9E\%E5\%88\%BA\%E6\%BF\%80\%E8\%96\%AC\%E7\%B2\%BE\%E7\%A5\%9E\%E7\%97\%85}{Wikipedia} with \textbf{18\% experiencing psychosis lasting beyond one month} \href{https://pmc.ncbi.nlm.nih.gov/articles/PMC4766712/}{PubMed Central} (Yui K, Goto K, Ikemoto S, Ishiguro T. ``Methamphetamine psychosis: spontaneous recurrence of paranoid-hallucinatory states and monoamine neurotransmitter function.'' \emph{Journal of Clinical Psychopharmacology} 1997; 17: 34--43). Tomiyama's comparative study of 11 chronic methamphetamine psychosis patients versus 11 matched chronic schizophrenia patients found that while positive symptoms were nearly identical, methamphetamine patients showed \textbf{less severe affective flattening and poverty of thought} \href{https://pubmed.ncbi.nlm.nih.gov/2074612/}{PubMed} \href{https://cir.nii.ac.jp/crid/1050570022170482432}{CiNii Research} and could rapidly switch between withdrawn and engaged states --- a qualitative difference captured by the absence of the \emph{Praecox-Gefühl} (Rümke's ``praecox feeling'') characteristic of schizophrenia \href{https://cir.nii.ac.jp/crid/1050570022170482432}{CiNii Research} (Tomiyama G. ``Chronic schizophrenia-like states in methamphetamine psychosis.'' \emph{Japanese Journal of Psychiatry and Neurology} 1990; 44(3): 531--539).

\subsection{What Tatetsu and colleagues predicted about recovery}\label{12_philopon-what-tatetsu-and-colleagues-predicted-about-recovery}

Tatetsu Seijun (立津政順, 1915--1999), the Okinawa-born psychiatrist who served at Matsuzawa Hospital from 1942 \href{https://ja.wikipedia.org/wiki/\%E7\%AB\%8B\%E6\%B4\%A5\%E6\%94\%BF\%E9\%A0\%86}{Wikipedia} through the entire first epidemic, was the founding figure of methamphetamine psychosis research worldwide. His 1956 monograph \emph{Kakuseizai Chūdoku} (覚醒剤中毒, ``Stimulant Intoxication''), co-authored with Goto Akira and Fujiwara Toshiko and published by Igaku-Shoin (医学書院), \href{https://journals.sagepub.com/doi/10.1177/009145090803500410}{Sage Journals} remains the definitive clinical text from the epidemic era --- 311 pages documenting symptomatology, clinical course, and institutional data from Matsuzawa Hospital. Tatetsu reported that psychosis generally did not \emph{progress} after drug cessation but noted that patients with hereditary predisposition to psychiatric disorders and pre-existing personality abnormalities were overrepresented among chronic cases. \href{https://ja.wikipedia.org/wiki/\%E7\%B2\%BE\%E7\%A5\%9E\%E5\%88\%BA\%E6\%BF\%80\%E8\%96\%AC\%E7\%B2\%BE\%E7\%A5\%9E\%E7\%97\%85}{Wikipedia} He estimated that \textbf{5--6\% of psychotic cases became indistinguishable from schizophrenia or bipolar disorder} \href{https://ja.wikipedia.org/wiki/\%E7\%B2\%BE\%E7\%A5\%9E\%E5\%88\%BA\%E6\%BF\%80\%E8\%96\%AC\%E7\%B2\%BE\%E7\%A5\%9E\%E7\%97\%85}{Wikipedia} --- a prognostic finding that proved prescient.

The concept of \textbf{reverse tolerance} (逆耐性, \emph{gyaku taisei}), developed experimentally by Utena Hiroshi (臺弘, 1913--2014) at Tokyo University and clinically elaborated by Sato Mitsumoto, became central to prognostic thinking. Unlike classical pharmacological tolerance where increasing doses are needed for the same effect, chronic methamphetamine exposure produces \textbf{sensitization} --- a lasting state in which progressively smaller doses, or eventually psychological stress alone, can trigger full psychotic relapse. \href{https://cir.nii.ac.jp/crid/1050001338388511744}{CiNii Research} \href{https://nyaspubs.onlinelibrary.wiley.com/doi/abs/10.1196/annals.1316.035}{Annals of the New York Academy of Sciences} Sato's landmark 1983 study \href{https://pubmed.ncbi.nlm.nih.gov/1553491/}{PubMed} examined 21 patients selected from 111 with chronic methamphetamine psychosis across 10 mental hospitals: \textbf{16 of 21 who re-used methamphetamine after long-term abstinence experienced acute paranoid psychosis nearly identical to their initial episode}. \href{https://pubmed.ncbi.nlm.nih.gov/1553491/}{PubMed} Four relapsed after a single dose smaller than what originally caused psychosis, and \textbf{one patient relapsed without any drug re-use at all}, \href{https://pubmed.ncbi.nlm.nih.gov/6860719/}{PubMed} triggered solely by psychological stress (Sato M, Chen CC, Akiyama K, Otsuki S. ``Acute exacerbation of paranoid psychotic state after long-term abstinence in patients with previous methamphetamine psychosis.'' \emph{Biological Psychiatry} 1983; 18(4): 429--440). This finding was later replicated in Akiyama's study of 80 incarcerated women: \href{https://pubmed.ncbi.nlm.nih.gov/21477052/}{PubMed} \textbf{25\% experienced spontaneous psychosis recurrence} during long-term abstinence, and those with spontaneous recurrence actually had \textbf{longer abstinence periods} (mean 17.5 months), suggesting the vulnerability does not diminish with time \href{https://pubmed.ncbi.nlm.nih.gov/21477052/}{PubMed} (Akiyama K, Saito A, Shimoda K. ``Chronic methamphetamine psychosis after long-term abstinence in Japanese incarcerated patients.'' \emph{American Journal on Addictions} 2011; 20(3): 240--249).

Prognostic factors identified during and after the epidemic include duration and dosage of use (persistent psychosis associated with ≥5 years of use history), \href{https://pmc.ncbi.nlm.nih.gov/articles/PMC5027896/}{PubMed Central} route of administration (injection producing more severe outcomes, though Matsumoto found smoking was not meaningfully safer for psychosis risk), \href{https://pubmed.ncbi.nlm.nih.gov/12133119/}{PubMed} comorbid psychiatric conditions, premorbid personality, and --- more recently --- \textbf{dopamine D2 receptor gene variants} (DRD2). Ujike and colleagues demonstrated that DRD2 -141C Del positive genotype conferred a 3.6-fold increased risk for rapid-onset psychosis, while DRD2 A1/A1 homozygosity was protective against poor prognosis \href{https://link.springer.com/chapter/10.1007/978-94-007-0837-2_15}{Springer} \href{https://pubmed.ncbi.nlm.nih.gov/19275926/?dopt=Abstract}{PubMed} (Ujike H, et al.~``Genetic variants of D2 but not D3 or D4 dopamine receptor gene are associated with rapid onset and poor prognosis of methamphetamine psychosis.'' \emph{Progress in Neuro-Psychopharmacology and Biological Psychiatry} 2009; 33(4): 625--629). The definition of ``recovery'' itself was contested: clinicians distinguished \textbf{symptom remission} (disappearance of hallucinations and delusions), \textbf{functional recovery} (return to employment and social roles), and \textbf{vulnerability status} (persistence of sensitization even in the absence of overt symptoms). By the latter criterion, no chronic user ever fully recovered.

\subsection{Thirty years later, nine patients still bore the scars}\label{12_philopon-thirty-years-later-nine-patients-still-bore-the-scars}

The long-term fate of chronic users from the first epidemic represents the most important and most poorly documented aspect of this history. The evidence comes from three tiers of declining reliability: a handful of follow-up studies (strong evidence), institutional and police records (moderate evidence), and extrapolation from clinical knowledge and modern analogy (weak evidence).

\textbf{The follow-up studies} constitute the strongest evidence. Tatetsu's 1997 study --- published when the author was 82 years old, near the end of his career --- followed \textbf{nine of his original Matsuzawa Hospital patients for 30 years} after their initial methamphetamine intoxication. The study's keywords reveal the findings starkly: \textbf{``schizophrenic state''} and \textbf{``remarkable disturbance of spontaneity''} persisted three decades after last drug use in at least some of these patients \href{https://webview.isho.jp/journal/detail/abs/10.11477/mf.1405904411}{Isho} (Tatetsu S, Hanawa S, Hamamoto J. ``覚醒剤中毒者9人の30年後.'' \emph{Seishin Igaku} 1997; 39(10): 1035--1043). Teraoka's earlier study tracked \textbf{114 former users 8--12 years} after chronic intoxication \href{https://psychiatryonline.org/doi/pdf/10.1176/appi.ajp.2009.08111695?download=true}{Psychiatry Online} and found \textbf{28\% exhibiting schizophrenia-like symptoms} at follow-up \href{https://psychiatryonline.org/doi/10.1176/appi.ajp.2009.08111695}{Psychiatry Online} --- nearly three times Tatetsu's original 1956 estimate of 5--6\% for chronic psychosis, likely reflecting the accumulation of residual and relapse cases over time (Teraoka A. ``Prognostic and criminological studies on the mental disorder caused by chronic methamphetamine intoxication.'' \emph{Seishin Shinkeigaku Zasshi} 1967; 69: 597--619).

\textbf{Institutional records} provide the most concrete quantitative data on the worst outcomes. Goto Akira's data from Matsuzawa Hospital, likely the source of figures later cited in the 1980 Police White Paper, documented that of \textbf{146 methamphetamine addicts admitted between 1947 and 1957}, 23 remained hospitalized as of August 1959 --- averaging \textbf{6 years and 10 months after cessation} --- and 15 of these 23 still exhibited recurrent psychotic symptoms. By May 1963, \textbf{17 patients remained hospitalized}, averaging \textbf{12 years and 2 months post-cessation}. Of these 17, \textbf{11 were confined to closed wards} and \textbf{5 were entirely unable to engage in any occupational activity} \href{https://www.npa.go.jp/hakusyo/s55/s550200.html}{National Police Agency} (Goto T. ``Clinical pictures shown by long hospitalized cases of chronic methamphetamine psychosis: a comparative study with schizophrenia.'' \emph{Psychiatria et Neurologia Japonica} 1960; 62: 163--176; data also cited in NPA 1980 Police White Paper). This represents approximately \textbf{12\% of admitted patients} remaining institutionalized over a decade after their last dose --- a figure that should be understood as a lower bound, since patients who died, were transferred, or were released in deteriorated condition are not counted.

Later clinical data, predominantly from the second epidemic (1970s--1980s), fills in further detail. Iwanami and colleagues documented that \href{https://pubmed.ncbi.nlm.nih.gov/8085475/?dopt=Abstract}{PubMed} \href{https://pubmed.ncbi.nlm.nih.gov/8085475/}{PubMed} \textbf{15\% of 104 patients} admitted to Matsuzawa Hospital (1988--1991) required more than three months of hospitalization despite abstinence and antipsychotic medication, \href{https://onlinelibrary.wiley.com/doi/abs/10.1111/j.1600-0447.1994.tb01541.x}{Wiley Online Library} \href{https://pubmed.ncbi.nlm.nih.gov/8085475/}{PubMed} and \textbf{28\% had symptoms persisting beyond six months} of abstinence \href{https://psychiatryonline.org/doi/10.1176/appi.ajp.2009.08111695}{Psychiatry Online} (Iwanami A, et al.~``Patients with methamphetamine psychosis admitted to a psychiatric hospital in Japan.'' \emph{Acta Psychiatrica Scandinavica} 1994; 89: 428--432). Akiyama's study of 32 incarcerated women found psychosis persisting \textbf{5 to 31 months after last injection} in 31 of 32 patients, \href{https://pmc.ncbi.nlm.nih.gov/articles/PMC7004251/}{PubMed Central} with symptoms including auditory hallucinations (91\%), persecutory delusions (91\%), depressive mood (91\%), and suicidal ideation (69\%) \href{https://pubmed.ncbi.nlm.nih.gov/21477052/}{PubMed} \href{https://pubmed.ncbi.nlm.nih.gov/17105910/}{PubMed} (Akiyama K. ``Longitudinal clinical course following pharmacological treatment of methamphetamine psychosis which persists after long-term abstinence.'' \emph{Annals of the New York Academy of Sciences} 2006; 1074: 125--134). Sato reported in 1992 that \textbf{nearly 50\% of second-epidemic hospital admissions for methamphetamine psychosis were readmissions}, \href{https://webview.isho.jp/journal/detail/abs/10.11477/mf.1405904147}{Isho} with the most extreme cases readmitted more than 10 times \href{https://pmc.ncbi.nlm.nih.gov/articles/PMC5027896/}{PubMed Central} --- evidence of a revolving door driven by the sensitization mechanism.

\textbf{Mortality data is almost entirely absent.} No formal cohort mortality study tracking first-epidemic users was ever conducted. Suicide was clearly a significant risk: Akiyama found 29 of 80 incarcerated female patients (36\%) had attempted suicide, with premorbid psychiatric disorders conferring an 8.6-fold increased risk. \href{https://pubmed.ncbi.nlm.nih.gov/21477052/}{PubMed} Cerebrovascular events associated with chronic methamphetamine use have been documented in modern case reports. \href{https://www.jstage.jst.go.jp/article/jstroke/advpub/0/advpub_10992/_article/-char/ja/}{JST} But no systematic data on life expectancy, cause-specific mortality, or overall survival exists for the 1950s cohort --- \textbf{the single largest gap in the evidence base}.

\textbf{Regarding cognitive and motivational recovery}, the picture is more nuanced than early prognostic pessimism suggested. Modern neuroimaging studies by Volkow and colleagues showed that dopamine transporter (DAT) levels, reduced by approximately 24\% in active users (roughly equivalent to 40 years of neurological aging), \textbf{recovered to near-normal levels after ≥9 months of abstinence} --- though motor and memory impairments persisted despite DAT recovery. \href{https://www.ehd.org/health_meth_6.php}{EHD} Iudicello and colleagues conducted the first controlled longitudinal study of longer-term cognitive recovery, finding that \textbf{abstinent users showed comparable global neuropsychological performance to healthy controls at one year}, with cognitively impaired users at baseline showing disproportionately large improvement \href{https://pubmed.ncbi.nlm.nih.gov/20198527/}{PubMed} \href{https://pmc.ncbi.nlm.nih.gov/articles/PMC2911490/}{PubMed Central} (effect size d = 1.73 for processing speed) (Iudicello JE, et al.~\emph{Journal of Clinical and Experimental Neuropsychology} 2010). However, \textbf{striatal metabolic deficits --- particularly in the nucleus accumbens --- persisted even after 12--17 months of abstinence}, which Volkow proposed as the neurobiological substrate for the persistent amotivation and anhedonia \href{https://pubmed.ncbi.nlm.nih.gov/14754772/}{PubMed} that Tatetsu had described clinically decades earlier. The Japanese clinical literature on residual syndrome (残遺症候群) documented anxiety, agitation, decreased motivation, personality changes, emotional disturbance, and fatigue \href{https://nyaspubs.onlinelibrary.wiley.com/doi/10.1111/j.1749-6632.2000.tb05178.x}{Annals of the New York Academy of Sciences} that proved \textbf{resistant to pharmacological treatment} and maintained heightened relapse vulnerability even in the absence of overt psychotic symptoms. \href{https://utu-yobo.com/column/9800}{Utu-yobo} \href{https://cir.nii.ac.jp/crid/1390006817998290944}{CiNii Research}

\textbf{Incarceration and recidivism} are better documented statistically, though primarily from the second and third epidemics. Modern longitudinal data from Hazama and Katsuta, tracking 1,807 paroled stimulant offenders over 10 years, found a \textbf{47.5\% drug-related recidivism rate}, \href{https://link.springer.com/article/10.1007/s11417-019-09299-8}{Springer} with stimulant offenders re-imprisoned an average of 2.9 times compared to 1.8 for violent offenders \href{https://www.tandfonline.com/doi/full/10.1080/23311886.2018.1489458}{Taylor \& Francis Online} (Hazama K, Katsuta S. ``Factors associated with drug-related recidivism among paroled amphetamine-type stimulant users in Japan.'' \emph{Asian Journal of Criminology} 2020). During the first epidemic, over 55,000 arrests occurred in the peak year alone, and approximately half of arrestees were stimulant-addicted. \href{https://www.unodc.org/unodc/en/data-and-analysis/bulletin/bulletin_1989-01-01_1_page007.html}{UNODC +2} Japan's zero-tolerance drug policy, rooted in the Hiropon-era campaigns, created enduring stigma that made social reintegration extraordinarily difficult: drug use was framed as ``innately selfish, socially and morally polluting,'' \href{https://www.vice.com/en/article/how-stigma-created-japans-hidden-drug-problem/}{VICE} and the DARC (Drug Addiction Rehabilitation Center) network reports that \textbf{60\% of recovering users express severe anxiety about finding employment}, \href{https://pubmed.ncbi.nlm.nih.gov/24946393/}{PubMed} while family rejection --- including being disowned --- is common. \href{https://pmc.ncbi.nlm.nih.gov/articles/PMC11527305/}{PubMed Central}

\textbf{What happened to most of the two million?} The dramatic arrest decline from 55,664 to 271 between 1954 and 1958, \href{https://www.unodc.org/unodc/en/data-and-analysis/bulletin/bulletin_1989-01-01_1_page007.html}{UNODC} \href{https://www.unodc.org/unodc/en/data-and-analysis/bulletin/bulletin_1989-01-01_1_page007.html}{unodc} combined with the absence of a persistent visible user population, strongly suggests that \textbf{the vast majority of users --- particularly casual and moderate users --- simply stopped} when supply was disrupted and penalties escalated. The 550,000 ``chronic users'' figure likely included a spectrum from heavy habitual users to moderately dependent individuals, and most appear to have been reabsorbed into Japanese society during the rapid economic growth of the late 1950s and 1960s. The truly chronic, heavily dependent subpopulation --- perhaps 10--15\% of chronic users, or 55,000--80,000 individuals based on Hayashi's and Teraoka's ratios --- faced grimmer trajectories: permanent psychosis vulnerability, periodic or chronic institutionalization, incarceration, social marginalization, or some combination thereof. Some likely resurfaced during the second epidemic beginning around 1970, \href{https://www.unodc.org/unodc/en/data-and-analysis/bulletin/bulletin_1989-01-01_1_page007.html}{UNODC} as Sato documented that formerly psychotic users could relapse after decades of abstinence upon re-exposure or stress. \href{https://www.jstage.jst.go.jp/article/pjmj/55/3/55_347/_article/-char/ja/}{JST +6}

\subsection{Annotated bibliography of key primary sources}\label{12_philopon-annotated-bibliography-of-key-primary-sources}

The following represents the most important sources for understanding the epidemic and its long-term consequences, ranked by their contribution to answering the central question of chronic user outcomes. Evidence quality is classified as \textbf{strong} (prospective/retrospective follow-up with defined cohort), \textbf{moderate} (clinical case series, hospital records, government statistics), or \textbf{weak} (individual case reports, extrapolation, opinion).

\textbf{Tatetsu S, Goto A, Fujiwara T.} \emph{Kakuseizai Chūdoku} (覚醒剤中毒). Tokyo: Igaku-Shoin, 1956. 311 pp.~Reprinted by Kimura Shoten, 1978. --- The foundational clinical monograph from the epidemic period. Documents symptomatology, clinical course, and institutional data from Matsuzawa Hospital. Estimated 5--6\% of psychotic cases became indistinguishable from schizophrenia. \href{https://ja.wikipedia.org/wiki/\%E7\%B2\%BE\%E7\%A5\%9E\%E5\%88\%BA\%E6\%BF\%80\%E8\%96\%AC\%E7\%B2\%BE\%E7\%A5\%9E\%E7\%97\%85}{Wikipedia} Provides the clinical taxonomy on which all subsequent research was built. \textbf{Strong evidence} (systematic clinical observation of large patient population). Available through Japanese academic libraries and NDL; no English translation exists.

\textbf{Tatetsu S, Hanawa S, Hamamoto J.} ``覚醒剤中毒者9人の30年後.'' \emph{Seishin Igaku} (精神医学) 1997; 39(10): 1035--1043. DOI: 10.11477/mf.1405904411. --- The only 30-year follow-up study of first-epidemic patients. Nine patients from Matsuzawa Hospital tracked for three decades. \href{https://webview.isho.jp/journal/detail/abs/10.11477/mf.1405904147}{Isho} Found persistent schizophrenic states and remarkable disturbance of spontaneity. \href{https://webview.isho.jp/journal/detail/abs/10.11477/mf.1405904411}{Isho} \textbf{Strong evidence} (longest follow-up study in the literature), though limited by very small sample size (n=9) and likely selection bias toward severe cases. Full text paywalled at isho.jp.

\textbf{Teraoka A.} ``Prognostic and criminological studies on the mental disorder caused by chronic methamphetamine intoxication.'' \emph{Seishin Shinkeigaku Zasshi} (精神神経学雑誌) 1967; 69: 597--619. --- The largest follow-up study of first-epidemic patients: 114 former users tracked 8--12 years. Found 28\% with schizophrenia-like symptoms at follow-up. \href{https://psychiatryonline.org/doi/10.1176/appi.ajp.2009.08111695}{Psychiatry Online} Combines psychiatric and criminological prognosis data. \textbf{Strong evidence} (largest cohort, substantial follow-up period). Published only in Japanese.

\textbf{Sato M, Chen CC, Akiyama K, Otsuki S.} ``Acute exacerbation of paranoid psychotic state after long-term abstinence in patients with previous methamphetamine psychosis.'' \emph{Biological Psychiatry} 1983; 18(4): 429--440. PMID: 6860719. --- Landmark paper establishing human methamphetamine sensitization. Twenty-one patients from 111 with chronic psychosis; 16 of 21 relapsed upon re-use after long abstinence, 4 after a single sub-threshold dose, 1 without any re-use. \href{https://pubmed.ncbi.nlm.nih.gov/6860719/}{PubMed} \textbf{Strong evidence} (systematic multi-hospital clinical study).

\textbf{Sato M, Numachi Y, Hamamura T.} ``Relapse of paranoid psychotic state in methamphetamine model of schizophrenia.'' \emph{Schizophrenia Bulletin} 1992; 18(1): 115--122. DOI: 10.1093/schbul/18.1.115. PMID: 1553491. --- Synthesizes over four decades of Japanese clinical research into the three-course classification (transient, prolonged, persistent) and the lasting-vulnerability concept. \href{https://pubmed.ncbi.nlm.nih.gov/1553491/}{nih +2} Reports 50\% readmission rate during the second epidemic. \textbf{Strong evidence} (comprehensive synthesis in a leading psychiatry journal). Freely accessible.

\textbf{Hayashi S.} ``Wake-amine addiction'' (ウェカミン中毒). \emph{Sōgō Igaku} (総合医学) 1955; 12: 656--661. --- Pre-neuroleptic era study of 74 patients. Two-thirds remitted within 20 days; approximately 10\% showed psychosis persisting beyond six months. \href{https://psychiatryonline.org/doi/10.1176/appi.ajp.2009.08111695}{Psychiatry Online} \textbf{Strong evidence} (original data from the epidemic period, untreated natural history). Published only in Japanese.

\textbf{Brill H, Hirose T.} ``The Rise and Fall of a Methamphetamine Epidemic: Japan 1945--55.'' \emph{Seminars in Psychiatry} 1969; 1(2): 179--194. \href{https://www.unodc.org/unodc/en/data-and-analysis/bulletin/bulletin_1980-01-01_1_page002.html}{UNODC} \href{https://www.cambridge.org/core/journals/journal-of-asian-studies/article/abs/methamphetamine-solution-drugs-and-the-reconstruction-of-nation-in-postwar-japan/0A0FCE7E74AE7D1ECE9A134114773CE6}{Cambridge Core} --- The foundational English-language account. Estimated 1--2 million involved in parenteral use. Described Japan's success as ``unique and without historical precedent.'' \href{https://www.unodc.org/unodc/en/data-and-analysis/bulletin/bulletin_1980-01-01_1_page002.html}{UNODC} \textbf{Moderate evidence} (secondary review, but authoritative and universally cited).

\textbf{Tamura M.} ``Japan: Stimulant Epidemics Past and Present.'' \emph{UNODC Bulletin on Narcotics} 1989; 41(1--2): 83--93. \href{https://www.unodc.org/unodc/en/data-and-analysis/bulletin/bulletin_1989-01-01_1_page007.html}{unodc} \href{https://www.semanticscholar.org/paper/The-forgotten-success-story:-Japan-and-the-problem-Edstr\%C3\%B6m/59c1dec41766481719ddc07bcc096c950ff9399d}{Semantic Scholar} Available at: https://www.unodc.org/unodc/en/data-and-analysis/bulletin/bulletin\_1989-01-01\_1\_page007.html --- The most comprehensive single English-language source with detailed statistics from the National Research Institute of Police Science. \href{https://www.unodc.org/unodc/en/data-and-analysis/bulletin/bulletin_1989-01-01_1_page007.html}{UNODC} Contains arrest data, seizure quantities, user estimates, and demographic breakdowns. \href{https://www.unodc.org/unodc/en/data-and-analysis/bulletin/bulletin_1989-01-01_1_page007.html}{UNODC} \href{https://www.unodc.org/unodc/en/data-and-analysis/bulletin/bulletin_1989-01-01_1_page007.html}{unodc} \textbf{Strong evidence} (official data from police research institute). Freely available online.

\textbf{Akiyama K, Saito A, Shimoda K.} ``Chronic methamphetamine psychosis after long-term abstinence in Japanese incarcerated patients.'' \emph{American Journal on Addictions} 2011; 20(3): 240--249. PMID: 21477052. --- Eighty incarcerated women with psychosis during long-term abstinence. \href{https://pubmed.ncbi.nlm.nih.gov/21477052/}{PubMed} Twenty-five percent experienced spontaneous psychotic recurrence; those with longer abstinence periods paradoxically had higher recurrence. \href{https://pubmed.ncbi.nlm.nih.gov/21477052/}{PubMed} \textbf{Strong evidence} (systematic clinical study with substantial sample size).

\textbf{Kingsberg M.} ``Methamphetamine Solution: Drugs and the Reconstruction of Nation in Postwar Japan.'' \emph{Journal of Asian Studies} 2013; 72(1): 141--162. \href{https://www.cambridge.org/core/journals/journal-of-asian-studies/article/abs/methamphetamine-solution-drugs-and-the-reconstruction-of-nation-in-postwar-japan/0A0FCE7E74AE7D1ECE9A134114773CE6}{Cambridge Core} \href{https://read.dukeupress.edu/journal-of-asian-studies/article-abstract/72/1/141/325329/Methamphetamine-Solution-Drugs-and-the}{Duke University Press} DOI: 10.1017/S0021911812001787. --- The most important recent English-language historical analysis. \href{https://www.cambridge.org/core/journals/journal-of-asian-studies/article/abs/methamphetamine-solution-drugs-and-the-reconstruction-of-nation-in-postwar-japan/0A0FCE7E74AE7D1ECE9A134114773CE6}{Cambridge Core} Uses extensive Japanese archival sources including Diet proceedings and propaganda materials. Argues the epidemic's suppression was intertwined with national identity reconstruction. \href{https://www.cambridge.org/core/journals/journal-of-asian-studies/article/abs/methamphetamine-solution-drugs-and-the-reconstruction-of-nation-in-postwar-japan/0A0FCE7E74AE7D1ECE9A134114773CE6}{Cambridge Core} \href{https://www.researchgate.net/publication/282513182_The_forgotten_success_story_Japan_and_the_methamphetamine_problem}{ResearchGate} \textbf{Strong evidence} (archival historical research). See also Kingsberg's monograph \emph{Moral Nation: Modern Japan and Narcotics in Global History} (University of California Press, 2014).

\subsection{Conclusion: what the evidence tells us, and what it cannot}\label{12_philopon-conclusion-what-the-evidence-tells-us-and-what-it-cannot}

The post-WWII Japanese methamphetamine epidemic was the world's first mass stimulant crisis, \href{https://www.jstage.jst.go.jp/article/faruawpsj/57/2/57_139/_pdf/-char/ja}{JST +3} and the Japanese psychiatric community's response produced foundational knowledge about methamphetamine psychosis that remains central to the field. \href{https://psychiatryonline.org/doi/10.1176/appi.ajp.2009.08111695}{Psychiatry Online} The clinical trajectory that emerges from seven decades of research follows a consistent pattern: most users recovered rapidly once supply was eliminated, but a substantial minority --- conservatively \textbf{10--15\% of those who developed psychosis, and approximately 28\% of chronic users followed long-term} --- suffered lasting psychiatric consequences. The mechanism now understood as behavioral sensitization (reverse tolerance) created a permanent alteration in brain reward and stress circuitry, \href{https://pubmed.ncbi.nlm.nih.gov/15542728/?dopt=Abstract}{PubMed} meaning that even after decades of abstinence and apparent recovery, former chronic users remained \textbf{one stressor or one small dose away from full psychotic relapse}. \href{https://pubmed.ncbi.nlm.nih.gov/11085329/}{PubMed +4} This is not metaphor: it is documented clinical reality \href{https://pubmed.ncbi.nlm.nih.gov/15542728/?dopt=Abstract}{PubMed} from Sato's 1983 study \href{https://pubmed.ncbi.nlm.nih.gov/6860719/}{PubMed} through Akiyama's 2011 prison data.

The evidence base has three critical weaknesses that must be stated plainly. First, \textbf{no prospective cohort study ever tracked a representative sample of first-epidemic users through to natural death}. The follow-up studies that exist (Teraoka n=114, Tatetsu n=9) are small and likely biased toward severe cases who remained in contact with the psychiatric system. Second, \textbf{mortality data is essentially nonexistent} --- we cannot determine whether chronic 1950s users died younger, or from what causes. Third, \textbf{social outcome data is fragmentary}: employment, housing, family status, and community reintegration were documented anecdotally rather than systematically. The dominant narrative of ``epidemic resolved'' served Japan's national reconstruction project \href{https://www.cambridge.org/core/journals/journal-of-asian-studies/article/abs/methamphetamine-solution-drugs-and-the-reconstruction-of-nation-in-postwar-japan/0A0FCE7E74AE7D1ECE9A134114773CE6}{Cambridge Core} but obscured the individual fates of hundreds of thousands of people. What we can say with confidence is that for the chronic, heavily dependent subpopulation, the epidemic's end did not mean recovery --- it meant the beginning of a quieter, less visible form of suffering that, for some patients documented at Matsuzawa Hospital, continued for decades behind the walls of closed psychiatric wards.
